\section{Estado del arte}
\label{sec:arte}

Las guías de INSARAG dividen la búsqueda en cinco niveles; el biobot corresponde al nivel~4, la búsqueda
secundaria en todos los huecos donde alguien podría seguir vivo, después de que el perro marca una zona y
antes de mover losas~\cite{insarag}. La tabla~\ref{tab:tecnologias} compara las herramientas disponibles y
la tabla~\ref{tab:grupos}, los trabajos con insectos cíborg.

\begin{table}[!htb]
\centering\footnotesize
\caption{Tecnologías de búsqueda en estructuras colapsadas.}
\label{tab:tecnologias}
\rowcolors{2}{SoftGray}{white}
\begin{tabularx}{\textwidth}{@{}>{\bfseries\raggedright\arraybackslash}p{2.5cm} >{\raggedright\arraybackslash}p{3.0cm} >{\raggedright\arraybackslash}p{3.8cm} >{\raggedright\arraybackslash}X@{}}
\toprule
\rowcolor{TableHead} Tecnología & Costo por unidad & Movilidad en rendijas & Límite principal \\
\midrule
Binomio canino & 50\,000 a 75\,000 USD~\cite{sdf} & Trabaja la superficie & Ciclos de 20 a 45~min; falla con fuego, lluvia o calor~\cite{murphy2004} \\
Cámara en poste & Sin dato & 2 a 5.5~m en línea casi recta & No dobla esquinas~\cite{murphy2021} \\
Radar FINDER & Sin dato & No entra: ve a través de 9~m de escombro & Ubica con $\pm1.5$~m; no dice por dónde llegar~\cite{jpl2015} \\
Robot de orugas & 50\,000 a 150\,000 USD~\cite{murphy2021} & Huecos de más de $6\times17$~cm, con cable & Cero personas vivas halladas hasta 2021 \\
Robot serpiente & Prototipo & 5~cm de diámetro, con cable & Una misión real, sin hallazgos~\cite{whitman2018} \\
Biobot & Sin dato publicado & Rendijas de unos 2~cm; trepa 1.5~cm & 3 a 6~cm/s; 2 a 4~h; un despliegue real, sin hallazgos~\cite{tranngoc2023,htx2025} \\
\bottomrule
\end{tabularx}
\end{table}

\begin{table}[!htb]
\centering\footnotesize
\caption{Trabajos de referencia con \Gp{} (salvo donde se indica).}
\label{tab:grupos}
\rowcolors{2}{SoftGray}{white}
\begin{tabularx}{\textwidth}{@{}>{\bfseries\raggedright\arraybackslash}p{3.2cm} >{\raggedright\arraybackslash}p{4.0cm} >{\raggedright\arraybackslash}X@{}}
\toprule
\rowcolor{TableHead} Grupo y año & Estímulo & Resultado \\
\midrule
NC State, 2012~\cite{latif2012} & Electrodo en la antena, 3~V & Control remoto; mochila de 4~g; 10~\% de los recorridos se completa \\
Bowdoin, 2015~\cite{erickson2015} & Antena y cercos, bifásico, 50~Hz & Responde bien el 56~\%; el pulso monofásico habitúa, el bifásico no \\
Texas A\&M, 2015~\cite{sanchez2015} & Ganglio protorácico (otra especie) & 60~\% de éxito en giros; se recupera con 10~min de descanso \\
NC State, 2017~\cite{dirafzoon2017} & Comandos al azar & Se mueve más (37 a 64~\% del tiempo) y deja la pared diez veces más \\
NTU, 2023~\cite{tranngoc2023} & Cercos, 40~Hz & 94~\% de éxito entre muros; detecta personas con matriz de $32\times32$ \\
NTU, 2023~\cite{lin2023} & Superficie, sin cirugía & 81~\% de éxito por estímulo en giro y 74~\% en avance \\
HIT, 2023~\cite{ma2023} & Antena y cerco, unipolar & La respuesta cae 57~\% tras diez estímulos \\
NTU y Osaka, 2025~\cite{bai2025} & Antenas y cercos & Enjambre de 20 insectos con la mitad de estímulos \\
NTU, 2025~\cite{lin2025} & Electrodo torácico & Ensamble robotizado en 68~s; mochila de 2.3~g \\
Osaka, 2025~\cite{refat2025} & Luz ultravioleta & Sin cirugía; sin habituación en 150 ensayos \\
UDLAP, 2026~\cite{gracidas2026} & Temperatura y luz & Actividad máxima cerca de 30~\textdegree C; la luz ultravioleta aumenta la locomoción \\
\bottomrule
\end{tabularx}
\end{table}

\textbf{Lo que nadie ha resuelto.} La revisión deja cinco vacíos que orientan el proyecto: (1) nadie cierra el
lazo solo con lo que va en la mochila, porque la navegación publicada depende de cámaras externas o de una
computadora~\cite{tranngoc2023,bai2025}; (2) cada mochila avisa por separado, y no encontramos un trabajo que
combine las detecciones de dos para ubicar a la persona ni que use a una como relevo de otra; (3) no hay
probabilidad de detección ni tiempo a la detección medidos para biobots; (4) la habituación se maneja caso por
caso, sin una regla que una mochila pueda aplicar, y (5) soltar un enjambre que reporta a un solo operador se
consideró «físicamente impráctico»~\cite{murphy2004}. La respuesta que se propone es que nadie pilotee: cada
mochila decide sola y solo habla cuando encuentra algo.
